%% SandboxBench --- paper skeleton.
%%
%% Status: this is a structure with its argument written out and its results
%% missing. Every number is a placeholder marked \needsdata, and the file will
%% not quietly pass as finished: see the draft banner on page 1. Remove
%% \sbdraft once the evaluation is real.
\documentclass[sigconf,nonacm]{acmart}

\settopmatter{printacmref=false}
\renewcommand\footnotetextcopyrightpermission[1]{}

\usepackage{booktabs}
\usepackage{xcolor}
\usepackage{amsmath}
\usepackage{graphicx}

%% \needsdata marks a claim that no experiment supports yet. It is deliberately
%% loud: a placeholder that looks like a result is how a draft turns into a
%% false claim.
\newcommand{\needsdata}[1]{\textcolor{red}{\textbf{[needs data: #1]}}}
%% \sbdraft prints the banner. Delete the line that uses it when the paper is
%% no longer a skeleton.
\newcommand{\sbdraft}{%
  \begin{center}\fcolorbox{red}{yellow!20}{%
  \parbox{0.92\columnwidth}{\small\textbf{Draft skeleton.} The argument and the
  benchmark design are written. No experiment has been run. Every quantitative
  claim is marked \needsdata{} and must not be cited.}}\end{center}}

\begin{document}

\title{SandboxBench: A Benchmark Suite for Cloud-based Code Execution Sandboxes}

\author{Jiwei Yuan}
\affiliation{%
  \institution{Independent}
  \country{}
}
\email{ji-weiyuan@outlook.com}

\begin{abstract}
Cloud providers now sell code execution sandboxes as a product: a service
creates an isolated environment on demand, runs commands in it, keeps its
filesystem across a suspend, and destroys it. AI coding agents and
reinforcement-learning training loops are the workloads driving that market,
and both of them use a sandbox in a way that function-as-a-service benchmarks
do not describe. A FaaS benchmark measures one invocation that is stateless by
construction. A sandbox session is interactive, holds state between commands,
is reset rather than re-invoked, and runs beside thousands of its neighbours.

We present SandboxBench, a benchmark suite for cloud-based code execution
sandboxes. It contributes a workload model for stateful interactive execution,
a benchmark specification of five performance dimensions crossed with two
execution scales plus four end-to-end application workloads, a
provider-agnostic measurement framework, and a cross-provider evaluation. The
framework reports a provider's missing capabilities as missing rather than as
zero scores, and separates the latency of a provider's command transport from
the speed of the sandbox it reaches, because those two are routinely conflated
and differ by an order of magnitude. \needsdata{the headline finding}
\end{abstract}

\keywords{sandboxes, benchmarking, cloud computing, AI agents, reinforcement learning}

\maketitle

\sbdraft

\section{Introduction}
\label{sec:introduction}

A code execution sandbox is now a product you buy by the API call. E2B,
Daytona, Runta, Modal and others expose the same four verbs --- create a
sandbox, run a command in it, suspend it, destroy it --- and the workloads
buying them are AI coding agents and reinforcement-learning training loops.
Both run untrusted, model-generated code, so both need an isolation boundary;
both run it repeatedly against evolving state, so neither fits the stateless
invocation that serverless platforms were built around.

The research community has a mature answer for serverless. SeBS established
what a serverless benchmark suite must do: pick representative applications,
build one framework that deploys and measures them across providers, and fix a
methodology for performance, cost and reproducibility~\cite{sebs2021}.
SeBS-Flow then showed what justifies a \emph{new} suite rather than more
benchmarks in the old one: a new execution model that the existing tests cannot
characterise, in its case function orchestration rather than single
functions~\cite{sebsflow2025}.

This paper argues that cloud code execution sandboxes are such a model, and
measures them.

\subsection*{Why existing benchmarks do not transfer}

Four properties separate a sandbox session from a function invocation, and each
one breaks a measurement that FaaS benchmarks rely on.

\begin{itemize}
  \item \textbf{It is interactive.} A session is a sequence of commands over a
  warm environment, not one call. The cost that matters is the per-command
  overhead on a sandbox that already exists, which a cold-start number does not
  predict.
  \item \textbf{It is stateful.} The filesystem survives between commands, and
  on some providers across a suspend. Correctness is therefore part of the
  measurement: a resume that silently cold-boots is fast and useless.
  \item \textbf{It is reset, not re-invoked.} A reinforcement-learning
  environment returns to a start state thousands of times per training run.
  Depending on the provider and the environment, that reset is an in-process
  operation, a sandbox recreate, or a snapshot restore. Reporting one number
  across the three compares nothing.
  \item \textbf{It is multi-tenant at scale.} A rollout fleet runs thousands of
  sandboxes concurrently. What a trainer waits on is not single-sandbox speed
  but sustainable provisioning throughput and the slowdown caused by
  neighbours.
\end{itemize}

\subsection*{Contributions}

\begin{description}
  \item[C1: A workload model.] We characterise stateful, interactive code
  execution: session length, command duration distribution, the split between
  environment preparation and useful work, reset frequency, state size, and
  arrival pattern (\S\ref{sec:background}). \needsdata{the characterisation
  itself; without it the benchmark design is an engineering guess}
  \item[C2: A benchmark specification.] Five performance dimensions ---
  lifecycle, execution, compute, storage, network --- crossed with two
  execution scales, single and cluster, plus four end-to-end application
  workloads and a reinforcement-learning profile built from the same
  dimensions (\S\ref{sec:design}).
  \item[C3: A provider-agnostic framework.] An open-source Go implementation
  with a deliberately small provider contract, capability reporting that
  distinguishes \emph{unsupported} from \emph{slow}, and a timing method that
  separates a provider's command transport from the sandbox behind it
  (\S\ref{sec:implementation}).
  \item[C4: A cross-provider evaluation.] Measurements across three hosted
  providers answering four research questions (\S\ref{sec:evaluation}).
  \needsdata{all of it}
\end{description}

The framework and every configuration file in this paper are at
\url{https://github.com/jiweiyuan/sandboxbench}.

\section{Background and Motivation}
\label{sec:background}

\subsection{From function invocation to sandbox session}

A function-as-a-service platform offers one primitive: invoke a function,
return a result. Its benchmark suites follow that shape. SeBS measures
invocation latency, cold-start cost and resource use across six application
classes~\cite{sebs2021}; FunctionBench provides micro and application
workloads for the same model~\cite{functionbench2019}; vHive and the snapshot
line of work attack the cold start directly~\cite{vhive2021,catalyzer2020}.

A sandbox platform offers five: create, exec, suspend, resume, destroy. The
execution model is
\[
  \text{create} \rightarrow (\text{exec})^{n} \rightarrow
  [\text{suspend} \rightarrow \text{resume} \rightarrow (\text{exec})^{m}]
  \rightarrow \text{destroy},
\]
with $n$ and $m$ in the hundreds for an agent session and the bracketed
segment repeated for a long-lived one. Three measurement consequences follow.

First, cold start stops being the headline. It is paid once per session and
amortised over $n$ commands, so for large $n$ the per-command overhead
dominates the user-visible total. A ranking built on create latency alone can
inverse the ranking built on session completion time. \needsdata{whether it
does in practice --- this is RQ1}

Second, the exec path becomes a first-class object of study. It is not part of
the FaaS model at all, and providers have chosen materially different
mechanisms for it (\S\ref{sec:implementation}).

Third, correctness joins the measurement. Suspend and resume make the
filesystem a thing that can be silently lost, so a lifecycle benchmark that
reports only latency cannot distinguish a fast resume from a fresh boot.

\subsection{What the workloads actually do}
\label{sec:background:characterisation}

Benchmark design without workload data is guesswork dressed as rigour. Before
fixing the suite we therefore need answers to:

\begin{enumerate}
  \item How many commands does one sandbox session run, and how is that
  distributed?
  \item How are individual command durations distributed? Where is the tail?
  \item What fraction of session wall-clock goes to environment preparation,
  to useful execution, and to waiting for the next interaction?
  \item How often is a reinforcement-learning environment reset, and by which
  mechanism?
  \item How large is the state that must survive a suspend?
  \item What is the concurrency and the request arrival pattern of a real
  rollout fleet?
\end{enumerate}

\needsdata{answers from agent execution traces, public RL environments, and
aggregated provider metrics. Where a trace cannot be obtained, a synthetic
workload is used and labelled as synthetic, with its representativeness stated
as a limitation rather than assumed.}

\section{Benchmark Design}
\label{sec:design}

We adopt the four design principles SeBS states --- relevance, usability, a
scientific evaluation methodology, and extensibility~\cite{sebs2021} --- and
add one the sandbox setting forces on us: \emph{capability honesty}. Providers
do not offer the same operations, and a suite that scores a missing operation
as zero produces a ranking that rewards having fewer features.

\subsection{Five dimensions, two scales}

Table~\ref{tab:dimensions} is the specification. The rows are dimensions: what
system capability is being measured. The columns are scales: one sandbox, or
many at once.

\begin{table}[t]
\caption{The benchmark specification. Scalability is not a sixth row: it is the
cluster column of the other five.}
\label{tab:dimensions}
\small
\begin{tabular}{@{}lll@{}}
\toprule
Dimension & Single sandbox & Cluster \\
\midrule
Lifecycle & create-to-first-exec & ready throughput, burst p99 \\
Execution & per-command latency & concurrent command throughput \\
Compute   & build completion time & concurrent build slowdown \\
Storage   & metadata ops/s, IOPS & concurrent I/O slowdown \\
Network   & egress RTT, throughput & concurrent network slowdown \\
\bottomrule
\end{tabular}
\end{table}

Two choices in that table need defending.

\paragraph{Scalability is not a dimension.} An earlier version of this design
listed scalability beside lifecycle, compute and storage. That was wrong.
Scalability is not an independent resource; it is what happens to the other
five under load. Keeping it as a peer row would score the same underlying
behaviour twice --- once as ``storage'' and again as ``scalability'' ---
and any composite ranking built on the result would double-count it.

\paragraph{Orthogonal does not mean statistically independent.} The five
dimensions have distinct \emph{measurement targets}, not uncorrelated outputs.
A build time depends on CPU, filesystem and memory together. What the design
requires is weaker and achievable: each dimension measures a different thing,
each microbenchmark isolates one dominant bottleneck, each result answers a
question no other result answers, and no single underlying behaviour is scored
more than once.

\subsection{Microbenchmarks}

Eight microbenchmarks populate the table. Each exists to explain a number in
the application results, not to be a ranking of its own.

\begin{description}
  \item[B01 Cold create.] Create a sandbox and run the first command.
  Measured to the completion of that command, not to the control plane
  reporting ``running'': a provider that returns early has moved its startup
  cost into the user's first command rather than removed it.
  \item[B02 Resume.] Write a sentinel, suspend, resume, run a command, read
  the sentinel back. Reports resume-to-ready latency \emph{and} a state
  correctness rate.
  \item[B03 Lifecycle churn.] Create and destroy continuously. Reports
  latency drift and cleanup failures; a provider that leaks sandboxes after a
  burst is charging for them.
  \item[B04 Short command.] A no-op command on a live sandbox, repeated. This
  measures the exec transport, which for an interactive session is the cost
  that is paid $n$ times.
  \item[B06 CPU build.] Compile and run a fixed C program written by the
  benchmark itself, so the measurement never contains a package download and
  never shifts when an upstream toolchain does.
  \item[B09 Small files.] Create, stat and delete 5{,}000 small files. An
  unpacked dependency tree is a metadata workload, not a bandwidth one, so
  this is the storage number that predicts \texttt{npm ci}.
  \item[B12 Egress.] Fetch a fixed endpoint from inside the sandbox, timed by
  \texttt{curl} so the exec transport is excluded. Both the endpoint and the
  client location are part of the result; change either and the number stops
  comparing.
  \item[B14--B16 Cluster.] Burst provisioning, sustained open-loop
  provisioning, and noisy-neighbour isolation (\S\ref{sec:design:cluster}).
\end{description}

Four candidate microbenchmarks were cut. Browser automation and data
processing were dropped because both re-measure compute, storage and network
without isolating anything; parallel execution and long-running execution were
dropped as benchmarks and kept as \emph{parameters}, since the same workload
run at higher concurrency or for longer is not a different benchmark. Model
inference and model training were cut because they would make the attached
accelerator, not the sandbox, the thing being ranked.

\subsection{Four application workloads}

Application workloads describe what a developer or an agent waits for. They
mix dimensions on purpose and are reported in their own table.

\begin{description}
  \item[A. Environment preparation.] \texttt{git clone} at a pinned commit,
  then a dependency install. Dominated by network and filesystem metadata.
  \item[B. Code build and execution.] Build, then test. Dominated by CPU,
  memory and disk.
  \item[C. Interactive development.] Start a dev server, wait for it to serve,
  change a file, wait for the reload. Exercises long-lived processes, file
  watching and port exposure.
  \item[D. Stateful continuation.] Write work, suspend, resume, verify. The
  only application workload with a correctness result as well as a latency.
\end{description}

No application workload calls a language model. Model latency and output
non-determinism would be added to a measurement of infrastructure, and would
dominate it.

\subsection{The reinforcement-learning profile}

The RL profile reuses the dimensions above rather than adding new ones: a
reset is lifecycle plus storage, a step is execution, a rollout is cluster
scale. It contributes three benchmarks.

\begin{description}
  \item[Environment step.] Action to observation on a live environment, with
  fixed step compute and a configurable observation size, so what the series
  measures is the sandbox's own scheduling and transport overhead.
  \item[Environment reset.] Run three times, once per mechanism ---
  in-process, sandbox recreate, snapshot restore --- and reported as three
  results. An environment reset is not necessarily a sandbox restart, and
  collapsing the three would compare incomparable operations.
  \item[Parallel rollouts.] $N$ environments concurrently, each running a
  fixed episode. Reports transitions per second and the fraction of episodes
  that completed intact, because a rollout that breaks partway costs the
  trainer the whole episode.
\end{description}

\subsection{Lifecycle profiles}

The same workload is driven under four lifecycle patterns, which is why
duration and concurrency are run parameters rather than separate benchmarks:
\emph{ephemeral} (create, run once, destroy), \emph{persistent} (create, run
many times, destroy), \emph{recoverable} (create, run, checkpoint, resume) and
\emph{pooled} (prewarm, acquire, run, release).

\subsection{Cluster experiments}
\label{sec:design:cluster}

Three cluster metrics carry the scale story.

\paragraph{Sustainable provisioning throughput.} Not creation requests
accepted per second, but sandboxes that reach their first command per second,
under a stated service-level objective:
\[
  T_{\mathrm{sustained}} =
  \frac{\text{sandboxes that reached first exec}}{\text{elapsed seconds}}
  \;\Big|\; p_{99} \le L,\; \text{success} \ge S.
\]
We report $L$ and $S$ with every value. A control plane that accepts a
thousand creates and queues them has not provisioned a thousand sandboxes.

\paragraph{Effective density.} Not how many sandboxes a host admits, but how
many run SLO-compliant work:
\[
  D_{\mathrm{effective}} =
  \frac{\text{SLO-compliant active sandboxes}}{\text{physical hosts}}.
\]
For hosted providers the denominator is unobservable, so this metric applies
only to the self-hosted infrastructure configuration
(\S\ref{sec:methodology}).

\paragraph{Performance isolation.} Fixed work measured alone and then beside
$N$ interferers:
\[
  \text{slowdown} = \frac{t_{\mathrm{contended}}}{t_{\mathrm{baseline}}} - 1 .
\]

\subsection{Result status}

Every benchmark against every provider ends in exactly one of four states, and
all four are printed in every table.

\begin{description}
  \item[completed] The benchmark ran; its numbers are reportable.
  \item[unsupported] The provider offers no such capability. This is a
  statement about features, not about speed.
  \item[failed] The provider supports the operation and the run broke. The
  reason is printed in full.
  \item[skipped] The configuration did not request it.
\end{description}

The distinction is not cosmetic. Collapsing \emph{unsupported} into a zero
score makes a provider look slow for a feature it never claimed, and makes the
resulting leaderboard reward a smaller feature set.

\section{Framework Implementation}
\label{sec:implementation}

SandboxBench is about 5{,}100 lines of Go, with a further 700 of tests. Go rather than Rust because the
framework is a distributed load generator and an adapter surface, not a
virtualisation engine: goroutines and contexts express open-loop and
closed-loop generation directly, and the cost of contributing a new provider
adapter matters more to the project than the runner's own peak performance.
Measurements that need precision run \emph{inside} the sandbox
(\S\ref{sec:impl:timing}), where the runner's garbage collector cannot reach
them.

\subsection{Three layers}

\begin{description}
  \item[Provider adapter] owns resource lifecycle and command execution. It
  knows nothing about what is being measured.
  \item[Workload] decides what to measure. It receives either a ready sandbox
  or, when it needs to control lifecycle itself, the provider.
  \item[Runner] decides how to drive it: repetitions, concurrency, arrival
  rate, duration, timeouts, teardown.
\end{description}

The provider contract is five operations on a sandbox --- \texttt{Exec},
\texttt{WriteFile}, \texttt{ReadFile}, \texttt{Suspend}, \texttt{Resume},
\texttt{Destroy} --- plus \texttt{Create} and \texttt{Capabilities} on the
provider. \texttt{Suspend} and \texttt{Resume} return a sentinel
\texttt{ErrUnsupported}, which the runner turns into the \emph{unsupported}
status rather than into a failure or a zero.

Capabilities are checked \emph{before} any sandbox is created, so an
unsupported benchmark costs nothing rather than wasting a create.

\subsection{Three exec transports}
\label{sec:impl:transports}

Implementing three adapters surfaced a finding that we did not expect and that
shapes the rest of the paper: the three providers execute commands over three
different mechanisms, with three different per-command cost structures
(Table~\ref{tab:transports}).

\begin{table}[t]
\caption{Command execution mechanisms. Compiled from each provider's published
API specification while writing the adapters.}
\label{tab:transports}
\small
\begin{tabular}{@{}llp{3.4cm}@{}}
\toprule
Provider & Transport & Per-command cost \\
\midrule
Daytona & REST            & one HTTP request to a per-sandbox toolbox proxy \\
E2B     & Connect-RPC     & a server-streaming RPC to \texttt{envd} inside the sandbox \\
Runta   & SSH             & a session on a multiplexed connection through an SSH gateway \\
\midrule
local   & process         & \texttt{os/exec} on the runner host (reference floor) \\
\bottomrule
\end{tabular}
\end{table}

Runta's public REST API has no execute endpoint at all; it covers lifecycle,
files and checkpoints, and commands go in over the SSH gateway. A suite that
reported a single ``command latency'' column across these three would mostly be
ranking transports. The transport is therefore recorded with every result and
printed in every capability table.

\subsection{Separating the transport from the sandbox}
\label{sec:impl:timing}

Every command is timed twice.

\begin{description}
  \item[Wall] is measured by the runner and includes one control-plane round
  trip.
  \item[In-box] is measured inside the sandbox by a shell wrapper that
  brackets the command and writes the elapsed nanoseconds to standard error,
  from where the runner recovers it.
\end{description}

Compute, storage and application workloads report in-box durations. Lifecycle
and the execution dimension report wall durations, because for those the round
trip is the thing under study. A sandbox whose shell cannot produce
nanosecond timestamps yields no in-box value, and the framework then falls back
to wall and records which source it used, rather than silently reporting a
zero.

\subsection{Load generation}

Cluster benchmarks generate load in two modes. Closed-loop holds a fixed
concurrency and starts new work as old work finishes. Open-loop sends requests
on a fixed-rate ticker without waiting for the previous one, so a slow provider
shows up as a growing latency tail rather than as a quietly reduced offered
load.

Two client-side confounders are controlled explicitly. The HTTP transport is
configured with a large connection pool, because with Go's default pool a
burst of a thousand creates queues on the client's own connections and the
benchmark then measures the runner. And HTTP 429 is surfaced as a typed error
rather than retried: a throughput ceiling that is really an account quota must
not be published as a throughput ceiling.

\subsection{Teardown}

Every sandbox is destroyed on a context derived with
\texttt{context.WithoutCancel}, so an interrupted or timed-out run still
releases what it allocated. Cleanup failures are counted and reported; a
provider that leaks sandboxes after a burst is billing for them.

\subsection{Reproducibility}

A run writes one JSON file containing the applied configuration, the client
platform and operator-supplied location, every summary, and every raw sample,
so any statistic in the paper can be recomputed from the artefact. Adapter
configuration blocks hold credentials, so only their \emph{key names} are
echoed. Bootstrap confidence intervals are seeded from the sample values
themselves, which makes a reported interval reproducible from the archived
samples alone. Intervals are suppressed below ten samples rather than
computed on data that cannot support them.

The same JSON renders to Markdown for reading and to a \texttt{booktabs}
tabular for this paper, so no number is retyped between the artefact and
Table~\ref{tab:results}.

\section{Experimental Methodology}
\label{sec:methodology}

\subsection{Two leaderboards that must not be merged}

\begin{description}
  \item[Provider benchmark.] Hosted services measured through their public
  APIs. The result includes the public network path, the control plane, the
  scheduler, the storage layer and any warm pool.
  \item[Infrastructure benchmark.] Firecracker, Cloud Hypervisor and gVisor
  deployed on identical hardware under our own
  control~\cite{firecracker2020,gvisor2019}.
\end{description}

The two are never ranked in one table. A hosted provider's CPU generation,
storage class, region and prewarming strategy are unknown and differ between
providers, so a hosted number and a controlled-hardware number do not measure
the same quantity.

\subsection{What every result publishes}

A sandbox benchmark measures a network path and an account configuration as
much as it measures a sandbox. Each result therefore carries: requested vCPU,
memory and disk; region or target; image or template identity; the client
machine and its network location; concurrency and arrival rate; warm or cold
state; whether persistence was enabled; the exec transport; and the wall-clock
window of the run. The framework records all of these automatically except the
client's network location, which the operator supplies.

\subsection{Statistics}

Each latency series is reported as $n$, median, $p_{90}$, $p_{95}$, $p_{99}$,
minimum and maximum, with a percentile bootstrap confidence interval on the
mean. We report medians rather than means as headline numbers: sandbox create
latency is long-tailed, and the mean of ten creates where one hit a cold pool
describes neither the warm case nor the cold one.

Repetitions are stated per benchmark. Runs are repeated on at least three
separate days, because provider performance varies with load that we cannot
observe, and a single day's measurement cannot be distinguished from that
day's conditions. \needsdata{the repetition schedule actually used}

\subsection{Cost}

Cost is reported per unit of completed work, not per hour:
\[
  C_{\mathrm{task}} = \frac{\text{total charge}}
  {\text{successfully completed tasks}},
\]
which charges startup overhead, execution time and failed retries to the work
they were spent on. For the reinforcement-learning profile the unit is a
valid rollout. \needsdata{prices, and the billing granularity each provider
actually applies}

\section{Evaluation}
\label{sec:evaluation}

\needsdata{this entire section. The structure below is the plan; the tables
are wired to the artefact generator and currently contain a smoke-test run
against the local reference adapter, which is not a provider comparison.}

We evaluate three hosted providers --- E2B, Daytona and Runta --- plus the
local reference adapter as a latency floor.

\subsection*{RQ1: Do serverless metrics predict sandbox performance?}

This is the question a reviewer asks first: why another suite? We test it
directly. For each provider we measure a FaaS-style baseline (cold start plus
one single-shot execution) and three sandbox-shaped workloads (100 sequential
short commands; install, edit, resume; reset followed by 100 steps), and ask
whether the provider ranking is stable across them and whether cold-start
latency correlates with session completion time.

If the rankings agree and the correlation is strong, the honest conclusion is
that existing metrics suffice and this suite is not needed. If they diverge,
SandboxBench has an experimental justification rather than an argued one.
\needsdata{which of the two it is}

\subsection*{RQ2: How do bottlenecks move with session shape?}

We sweep command count, interaction frequency and state size, and locate the
crossover where per-command overhead overtakes create latency in the
user-visible total.

\subsection*{RQ3: How does single-sandbox performance relate to cluster throughput?}

We compare each provider's single-sandbox dimension results against its burst
and sustained provisioning results, and test whether the fastest single
sandbox belongs to the provider with the highest sustainable throughput.

\subsection*{RQ4: How do lifecycle strategies affect rollout throughput and cost?}

We run the reinforcement-learning profile under all four lifecycle profiles
and compare transitions per second and cost per valid rollout across the three
reset mechanisms.

\subsection{Results}

\begin{table}[t]
\caption{Headline results. Generated by \texttt{sandboxbench report
--format=latex}; do not edit by hand. \needsdata{a real cross-provider run ---
the current contents are a local smoke test}}
\label{tab:results}
\small
% Placeholder. No benchmark run has produced this table yet.
%
% Regenerate it from a real run:
%   make -C paper table RESULT=results/<run>.json
%
% Do not hand-write numbers here. The point of generating it is that no value
% is retyped between the artefact and the paper.
\begin{tabular}{lrrr}
\toprule
Benchmark & E2B & Daytona & Runta \\
\midrule
\textit{no run yet} & --- & --- & --- \\
\bottomrule
\end{tabular}

\end{table}

\subsection{Threats to validity}

\begin{description}
  \item[Provider configuration is not controlled.] Hosted providers do not
  expose CPU generation, storage class or prewarming policy. A difference
  between two providers may be a hardware difference.
  \item[The client is a confounder.] Every number includes a public network
  path from one client location. We report the location and repeat from more
  than one. \needsdata{how many}
  \item[Synthetic workloads may not be representative.] Where
  \S\ref{sec:background:characterisation} could not be answered from real
  traces, the corresponding workload is synthetic and labelled as such.
  \item[Measurements age.] These are commercial services under active
  development. A result is a measurement of a date, and the date is in every
  result file.
\end{description}

\section{Related Work}
\label{sec:related}

\paragraph{Serverless benchmark suites.} SeBS defined the template this work
follows: a benchmark specification, a cross-provider platform, and an
evaluation methodology, under the principles of relevance, usability,
scientific evaluation and extensibility~\cite{sebs2021}. SeBS-Flow extended it
to function workflows and, more importantly for us, demonstrated what
justifies a new suite: an execution model the existing tests cannot
characterise~\cite{sebsflow2025}. FunctionBench contributed micro and
application workloads for FaaS~\cite{functionbench2019}, and ServiTrace argues
for fine-grained distributed tracing and end-to-end application measurement
rather than aggregate latency~\cite{servitrace2025}. All of them measure
stateless invocation. None measures an interactive stateful session, a
suspend-resume round trip, or an environment reset.

\paragraph{Sandbox and microVM infrastructure.} Firecracker established the
microVM as the isolation boundary for multi-tenant
code execution~\cite{firecracker2020}; the gVisor case study quantified what
syscall interception costs~\cite{gvisor2019}; SOCK and Catalyzer attacked
container and VM startup latency~\cite{sock2018,catalyzer2020}. This body of
work measures the runtime. We measure the service built on top of one, and
treat the runtime as unobservable for hosted providers.

\paragraph{Reinforcement-learning environment execution.} EnvPool shows that
environment throughput, not only policy computation, bounds RL training
throughput~\cite{envpool2022}. It optimises a local execution engine. The
question we ask is one layer down: how the same environment workload behaves
when its environments are cloud sandboxes from different providers.

\paragraph{Agent benchmarks.} SWE-bench and its successors evaluate what a
model can accomplish, using a sandbox as a fixture~\cite{swebench2024}. We
evaluate the fixture.

\paragraph{Industry measurement.} The Kubernetes agent-sandbox project
publishes a cluster performance assessment covering rapid burst, high-volume
ramp, steady-state churn, throughput and warm-pool
burst~\cite{k8sagentsandbox}; our cluster load profiles follow its shapes
directly. Vendor and community benchmarks of hosted sandbox providers also
exist; we treat them as a source of load patterns rather than as comparable
measurements, since their configuration and client location are generally not
published.

\section{Limitations and Conclusion}
\label{sec:conclusion}

SandboxBench measures CPU-based code execution sandboxes. GPU workloads,
model inference and model training are out of scope by design: including them
would make the attached accelerator the dominant variable and turn the suite
into a hardware comparison. Effective density applies only to the
self-hosted infrastructure configuration, because a hosted provider does not
expose a host count. And a hosted measurement is always a measurement of a
date; the suite records the date and expects to be re-run.

What we claim is a benchmark specification for an execution model that
serverless suites do not describe, an implementation that reports capability
honestly and separates a provider's command transport from the sandbox behind
it, and a cross-provider evaluation organised around whether existing metrics
predict sandbox behaviour at all. \needsdata{the findings}

Both the framework and the configuration files used here are open source at
\url{https://github.com/jiweiyuan/sandboxbench}.


\bibliographystyle{ACM-Reference-Format}
\bibliography{refs}

\end{document}
